\chapter{Auditoría del tercer run}\label{cap:auditoria}

\newcommand*{\veredicto}[2]{{\sffamily\bfseries\footnotesize\tikz[baseline=-0.6ex]\fill[#1] (0,0) circle[radius=2.6pt];\,\textcolor{#1}{#2}}}

Antes de abrir el pull request se auditó el tercer run: se contrastó cada frase del README con el código y con la
salida del kernel. Este capítulo resume esa auditoría, porque cambió lo que el trabajo puede afirmar.

\section{Qué se revisó y cómo}

La revisión tuvo dos partes. Un revisor independiente, el agente \texttt{ecc:mle-reviewer} en modo de solo lectura y
sin correr redes, leyó el código buscando fugas, comparaciones injustas y problemas de reproducibilidad. Después, la
auditoría recalculó desde las tablas por sujeto de la salida del kernel todo lo que el README afirma, y leyó
directamente los pesos publicados.

\begin{clave}
Las cifras cuadran. El \texttt{metricas.json} del repositorio es idéntico al del kernel, igual que los cuatro gráficos
(\cref{tab:traza}), y el revisor rehízo el rango esperado en CPU desde los parquet y obtuvo las mismas tablas. Tampoco
hay fuga: los cinco folds no se cruzan y ningún sujeto de prueba entra a un ajuste. Lo que la auditoría encontró son
conclusiones que dicen más de lo que muestran los datos y frases que no describen lo que hace el código.
\end{clave}

\section{Afirmación por afirmación}

La \cref{tab:afirmaciones} recorre las once afirmaciones del README que se revisaron, con su veredicto: se sostiene, falta
aclarar, es inexacta o está sobredicha.

\begingroup
\small
\renewcommand{\arraystretch}{1.25}
\begin{longtable}{@{}>{\raggedright\arraybackslash}p{4.9cm} >{\raggedright\arraybackslash}p{2.35cm}
  >{\raggedright\arraybackslash}p{7.05cm}@{}}
\caption[Las afirmaciones del README y su veredicto]{Las afirmaciones del README del tercer run y lo que muestran el
código y la salida del kernel.}\label{tab:afirmaciones}\\
\toprule
Afirmación & Veredicto & Por qué \\
\midrule
\endfirsthead
\multicolumn{3}{@{}l}{\footnotesize\sffamily\color{tinta3} Viene de la página anterior}\\
\toprule
Afirmación & Veredicto & Por qué \\
\midrule
\endhead
\midrule
\multicolumn{3}{@{}r}{\footnotesize\sffamily\color{tinta3} Sigue en la página siguiente}\\
\endfoot
\bottomrule
\endlastfoot
\enquote{Las cifras de este README salen del tercer run del kernel.} & \veredicto{serieC!80!black}{Se sostiene} &
\texttt{metricas.json} y los cuatro gráficos son idénticos a los del kernel (sha256). \\
\enquote{Los sujetos de prueba nunca entran a ningún ajuste.} & \veredicto{serieC!80!black}{Se sostiene} & El
ridge, el ancho calibrado, el piso de la desviación, la parada temprana y el ruido por variable salen solo de sujetos
de entrenamiento, y los folds no se cruzan. \\
\enquote{En el peso, el rango queda a la mitad y la detección sube de 10 a 26\pct.} &
\veredicto{medio}{Falta aclarar} & Es la comparación justa, pero ese 10\pct{} es el monitor de Tierra \emph{con el
ancho calibrado}, y la tabla lo llamaba \enquote{Tierra}. El de Tierra sin calibrar detecta \cifra{peso.tierra.rec}\pct,
con \cifra{peso.tierra.fareg}\pct{} de falsas alarmas por la regla (M2). \\
\enquote{Una alarma cuenta cuando 2 de 3 días seguidos quedan fuera, y esa misma regla se usa para las falsas
alarmas.} & \veredicto{serieB}{Inexacto} & La tabla usaba falsas alarmas diarias. Con la regla, el adaptado en la
presión diastólica da \cifra{pa_dia.adapt.fareg}\pct, no \cifra{pa_dia.adapt.fa}\pct{} (M1). \\
\enquote{El 5\pct{} del adaptado sale por construcción.} & \veredicto{medio}{Falta aclarar} & En promedio, sí. Pero
el ancho es uno para toda la población: en el peso, \cifra{fapeso.cero} de \cifra{fapeso.n} personas no tienen
ninguna falsa alarma y una llega a \cifra{fapeso.max}\pct{} (M4). \\
\enquote{El resultado fuerte sale del corrector: la línea base, lo aprendido de otros sujetos y los primeros 28
días.} & \veredicto{medio}{Falta aclarar} & Vale para el peso. Con los 28 días propios ya puestos, lo aprendido de
otros baja el error en más de \cifra{banda} solo en \cifra{aporte.ayuda}, y en el \cifra{aporte.empeora} lo sube (M3).
\\
\enquote{Con estos datos, el pre-entrenamiento en Tierra no aporta al pronóstico.} & \veredicto{alto}{Sobredicho} &
Se apoya en \enquote{solo la cabeza contra desde cero}. Con ajuste completo, Tierra → espacio gana
\cifra{par.todo.cero} \cifra{par.todo.cero.ic}, y el intervalo no cruza el cero. Tampoco se puede afirmar lo
contrario: una sola semilla y brazos no igualados (H1). \\
\enquote{Después se adapta al reposo ajustando solo la cabeza o todo el modelo.} & \veredicto{serieB}{Inexacto} &
\enquote{Solo la cabeza} también cambia el backbone: entre los pesos publicados cambian \cifra{pesos.cambian} de
\cifra{pesos.tensores} tensores, \cifra{pesos.bn} de BatchNorm (A1). \\
\enquote{NOVA aprende cómo es cada persona en la Tierra y se adapta a su nueva normalidad con los pocos días de la
misión.} & \veredicto{medio}{Falta aclarar} & Eso lo muestra el rango esperado. El pronóstico se adapta al reposo en
general: el sujeto de prueba nunca aporta sus propios días (M5). \\
\enquote{La mejora es real pero chica: cerca de 1,5\pct{} del error de persistencia.} &
\veredicto{serieB}{Cifra inexacta} & Son \cifra{par.cabeza.media7.puntos} puntos ajustando solo la cabeza y
\cifra{par.todo.media7.puntos} ajustando todo (L4). \\
\enquote{Con 11 a 21 días de pre-reposo y 42 a 90 en cama.} & \veredicto{serieB}{Cifra inexacta} & Contando los días
con pulso, son de \cifra{datos.chk.pulso.pre.min} a \cifra{datos.chk.pulso.pre.max} y de
\cifra{datos.chk.pulso.cama.min} a \cifra{datos.chk.pulso.cama.max} (L4). \\
\end{longtable}
\endgroup

\section{Hallazgos}

El revisor independiente encontró 0 críticos, 1 alto, 5 medios y 8 bajos; la auditoría agregó uno medio, el A1.
Cada hallazgo dice qué pasa y cómo se arregla.

\subsection{El hallazgo alto}\label{sec:h1}

\begin{hallazgo}{H1}{Alto}{alto}{\enquote{El pre-entrenamiento no aporta} sale de una comparación elegida}
\begin{description}
\item[Qué pasa] El README concluía \enquote{empate} comparando solo la cabeza contra el modelo desde cero,
  \cifra{par.cabeza.cero} \cifra{par.cabeza.cero.ic}. El ajuste completo, que es la mejor variante, no se comparaba con
  los modelos desde cero. Contra ellos gana: \cifra{par.todo.cero} \cifra{par.todo.cero.ic} y \cifra{par.todo.sinls}
  \cifra{par.todo.sinls.ic} sin LifeSnaps.
\item[Confusores] Tasa de aprendizaje distinta entre brazos; \enquote{solo la cabeza} cortada por el tope de épocas
  en \cifra{cabeza.cortados} de 5 folds; una sola semilla. Que los modelos desde cero sufran por su tasa de aprendizaje
  es una hipótesis: no se probó otra.
\item[Arreglo] Informar los pares del ajuste completo y escribir que la diferencia es de \cifra{par.todo.cero.r} o
  menos y que una semilla no permite separarla del ruido. En la corrida siguiente, igualar la tasa de aprendizaje y el
  presupuesto, subir el tope de la cabeza y usar varias semillas.
\end{description}
\end{hallazgo}

\subsection{Hallazgos medios}

\begin{hallazgo}{M1}{Medio}{medio}{La tabla usaba falsas alarmas diarias, no la regla de 2 de 3}
El texto decía que las falsas alarmas usan la misma regla que la detección, pero las columnas eran diarias. Con la
regla, el monitor de Tierra en la presión diastólica pasa de \cifra{pa_dia.tierra.fa} a
\cifra{pa_dia.tierra.fareg}\pct, y el adaptado de \cifra{pa_dia.adapt.fa} a \cifra{pa_dia.adapt.fareg}\pct.
\emph{Arreglo:} decir que la tabla es diaria y agregar la regla (\cref{tab:regla}).
\end{hallazgo}

\begin{hallazgo}{M2}{Medio}{medio}{La comparación de detección era contra el monitor de Tierra calibrado}
En el peso, el 10\pct{} es \modelo{tierra (ancho calibrado)}. El monitor de Tierra real detecta
\cifra{peso.tierra.rec}\pct, pero con \cifra{peso.tierra.fareg}\pct{} de falsas alarmas por la regla, así que esa
detección es ruido. En promedio sin agua, el adaptado detecta menos que el de Tierra sin calibrar
(\cifra{prom.adapt.rec} contra \cifra{prom.tierra.rec}\pct), y eso no se decía. \emph{Arreglo:} llamar a la columna
\enquote{Tierra con ancho calibrado → adaptado} y explicar por qué el monitor sin calibrar no sirve de referencia.
\end{hallazgo}

\begin{hallazgo}{M3}{Medio}{medio}{El corrector de otros sujetos suma solo en algunas variables}
Con los 28 días propios ya puestos, el corrector baja el error en más de \cifra{banda} solo en \cifra{aporte.ayuda}.
El README le daba el crédito en general; en la mayoría de las variables, la mejora viene de los días propios.
\emph{Arreglo:} mostrar \modelo{tierra+personal} junto al adaptado (\cref{tab:aporte}).
\end{hallazgo}

\begin{hallazgo}{M4}{Medio}{medio}{El 5\pct{} de falsas alarmas es de la población, no de cada persona}
En el peso, la falsa alarma diaria del adaptado por persona tiene mediana \cifra{fapeso.mediana}\pct, media
\cifra{fapeso.media}\pct{} y máximo \cifra{fapeso.max}\pct. Unas pocas personas alarman casi todos los días y la
mayoría nunca. \emph{Arreglo:} mostrar la distribución por persona (\cref{fig:fa-persona}) y probar un ancho por
persona.
\end{hallazgo}

\begin{hallazgo}{M5}{Medio}{medio}{El pronóstico no se adapta a cada persona}
La adaptación usa los días en cama de los sujetos de entrenamiento; el sujeto de prueba nunca ve los suyos. La
adaptación personal solo está en el rango esperado. \emph{Arreglo:} decirlo, o agregar un experimento con los primeros
$k$ días del sujeto de prueba.
\end{hallazgo}

\subsection{El hallazgo nuevo: BatchNorm se mueve}\label{sec:a1}

\begin{hallazgo}{A1}{Medio}{medio}{\enquote{Solo la cabeza} también re-estima BatchNorm del backbone}
\texttt{congelar\_backbone} solo apaga los gradientes, y el entrenamiento deja el modelo en modo de entrenamiento.
Entre \texttt{nova-tierra} y \texttt{nova-espacio} cambian \cifra{pesos.cambian} de \cifra{pesos.tensores} tensores:
\cifra{pesos.cabeza} de la cabeza y \cifra{pesos.bn} de BatchNorm (\cref{sec:pesos}). Puede explicar parte de lo que
\enquote{solo la cabeza} le gana a Tierra sin adaptar. \emph{Arreglo:} poner el backbone en modo de evaluación, o
describir el brazo como \enquote{cabeza + normalización}. Amplía el L1 del revisor.
\end{hallazgo}

\subsection{Hallazgos bajos}

\begin{hallazgo}{L1}{Bajo}{bajo}{Dropout activo en el backbone congelado}
El modo de entrenamiento también activa el dropout de las capas congeladas.
\end{hallazgo}
\begin{hallazgo}{L2}{Bajo}{bajo}{Reproducibilidad incompleta}
\texttt{chronos-forecasting>=2.0} se instala sin versión fija (en la corrida fue la 2.3.2), \texttt{metricas.json}
guarda el hash de los datos pero no el del código, y no se fija el determinismo de CUDA.
\end{hallazgo}
\begin{hallazgo}{L3}{Bajo}{bajo}{Ninguna aserción de que prueba y entrenamiento no se cruzan}
A mano está correcto, pero una aserción de una línea lo dejaría comprobado en cada corrida.
\end{hallazgo}
\begin{hallazgo}{L4}{Bajo}{bajo}{Cifras del README que no eran exactas}
El \enquote{cerca de 1,5\pct}, los días de pre-reposo y en cama, y que LifeSnaps tiene \cifra{datos.ls.personas}
personas pero al modelo base entran \cifra{datos.ls.entrenan}: 7 quedan para la parada temprana.
\end{hallazgo}
\begin{hallazgo}{L5}{Bajo}{bajo}{Las magnitudes de los escalones son supuestos}
Salvo la del pulso. El nivel de detección depende de ellas; lo robusto es el orden entre modelos. No se barrió la
magnitud.
\end{hallazgo}
\begin{hallazgo}{L6}{Bajo}{bajo}{Muchas comparaciones sin corrección}
Tres versiones del kernel, ocho variables y muchos pares, sin corrección por multiplicidad \parencite{benjamini1995}.
Conviene declararlo o quedarse con pares definidos de antemano.
\end{hallazgo}
\begin{hallazgo}{L7}{Bajo}{bajo}{Dos desplazamientos sin comentar}
El agua baja \cifra{desp.agua}\un{g/día} ($\pm$\cifra{desp.agua.ic}) y la presión diastólica sube
\cifra{desp.pa_dia}\un{mmHg} ($\pm$\cifra{desp.pa_dia.ic}); el README no los comentaba.
\end{hallazgo}
\begin{hallazgo}{L8}{Bajo}{bajo}{El cargador toma el primer número de cada celda}
Si una celda trae dos números, el segundo se pierde sin aviso. Conviene que el cargador avise cuando pase.
\end{hallazgo}

\section{Qué se corrigió en el pull request}

El pull request \#2 lleva el README y las preguntas abiertas corregidos con todo lo anterior:
\begin{itemize}
\item La columna de detección dice \enquote{Tierra con ancho calibrado → adaptado}, y el texto explica por qué el
  monitor sin calibrar no sirve de referencia.
\item La tabla dice que las falsas alarmas son diarias, y el texto da el rango con la regla de 2 de 3.
\item El 5\pct{} se presenta como un promedio de la población, con la distribución del peso por persona.
\item El crédito del corrector se limita al peso y el pulso.
\item \enquote{El pre-entrenamiento no aporta} se reemplazó por \enquote{todavía no se puede decir si aporta}, con los
  pares del ajuste completo y los confusores.
\item \enquote{Solo la cabeza} se describe como cabeza más normalización, con el conteo de tensores.
\item El pronóstico se describe como adaptado al reposo en general, no a cada persona.
\item Las cifras inexactas se corrigieron, y un párrafo dice cuáles cifras no están en \texttt{metricas.json} y de
  dónde salen.
\end{itemize}

Quedan para la corrida siguiente los cambios de código: igualar los brazos, congelar de verdad el backbone, fijar la
versión de Chronos, guardar el hash del código y agregar las aserciones. Una frase se corrigió después, el 15-11-2026:
el plan de la v4 en el README y la respuesta a la pregunta~2 decían que LifeSnaps \enquote{aporta solo el pulso}. Ahora
dicen lo correcto: LifeSnaps aporta sus ocho canales al pre-entrenamiento compartido y, de las variables de la v4,
solo el pulso también está en LifeSnaps.
